\section{System structure}  
\label{SystemModel}

After deriving the dynamics of all components in the network, the complete system can be drawn. In the appendix, \secref{systemdiagram}, the topology of the test system is described in detail. In the following, all the statements are based on the notations and relation system of the system figure in the appendix. 
\\
The way of modelling of the hydraulic system is analogous to an electric circuit. Most of the various hydraulic components can be represented as electronic equivalents and vice versa. In the block diagram of the system nodes are introduced which usually represent different potentials in an electronic system. In general, this is equivalent to hydraulic heads. In other words, nodes represent points in the system where pressure has different values. These nodes are introduced arbitrarily but taking into account that each individual component in the system has an effect on the pressure drop on their two corresponding endpoints. Therefore nodes are placed at all places where the modelling requires it. 
\\
In the network, volume flow rate is equivalent to current and the quantity of water has the same physical meaning as electric charge in an electric circuit. 
\\ 
Although nodes can be placed across all the endpoints of components in the system, some simplifications are carried out. There are two different types of simplification in the network. For a better transparency, these parts of the system are shown in \figref{fig:subsys_1} and \figref{fig:subsys_2}: 

%%picture of test system
%\begin{figure}[H]
%\centering
%\includegraphics[width=0.35\textwidth]{report/pictures/missingfigure}
%\caption{Simplification when the pump is not used}
%\label{fig:subsys_1}
%\end{figure}

\begin{figure}[H]
\centering
\begin{minipage}{0.5\textwidth}
  \centering
  \includegraphics[width=0.65\linewidth]{report/pictures/missingfigure}
  \captionof{figure}{$\omega = 0$}
  \label{fig:subsys_1}
\end{minipage}%
\begin{minipage}{0.5\textwidth}
  \centering
  \includegraphics[width=0.65\linewidth]{report/pictures/missingfigure}
  \captionof{figure}{$\omega \neq 0$}
  \label{fig:subsys_2}
\end{minipage}
\end{figure}

In \figref{fig:subsys_1} three components between $n_3$ and $n_{18}$ are shown. This is the part of the system where the WT connects. In this particular case the pump is turned off, however contributes to the pressure drop because its resistance is not zero. The same can be said for the valves, except that they are fully open but they modify the flow. Extra nodes are not introduced between the valves and the pump, instead they are seen as one component. This can be modelled by lumping the resistance of the valve, \eqref{CompleteValveModel}, into the model of the pipe, \eqref{PumpModel}, when the rotational speed is zero. Thus the following model yields for the case when $\omega = 0$:

\begin{equation}
  \Delta p = \Big(\frac{2}{k_{v100}^2} - a_{h2}\Big)|q| q 
  \label{omega_zero}
\end{equation}

And the following when $\omega \neq 0$

\begin{equation}
  \Delta p = \Big(\frac{2}{k_{v100}^2} - a_{h2}\Big)|q| q  + a_{h1} \omega_r q_i + a_{h0}{\omega_r}^2
  \label{omega_notzero}
\end{equation}

The case when $\omega = 0$ applies to the components between $n_4$-$n_8$, $n_3$-$n_{18}$  and to $n_5$-$n_{13}$. The case when $\omega \neq 0$ applies to the two main pumps, between $n_1$-$n_2$ and $n_1$-$n_7$. All these subsystems are modelled as described above. 